\section{Q-Learning: práctica de sistemas y monitoreo}

\subsection{Métricas de monitoreo y calidad de los datos}

\begin{knowledgebox}{Métricas de observación clave}
\begin{itemize}
    \item Cobertura del replay buffer: mide si las últimas $N$ transiciones cubren los pares estado-acción frecuentes.
    \item Distribución del TD error: monitorear la desviación estándar y el valor máximo del TD error permite detectar con anticipación una explosión numérica.
    \item Action distribution drift: verificar si la frecuencia de muestreo de la política de comportamiento se concentra en unas pocas acciones, para evitar que el entrenamiento derive hacia una política estrecha.
\end{itemize}
\end{knowledgebox}

Para lograr un despliegue estable, el equipo de ingeniería debe exponer estas métricas en el tablero de entrenamiento y definir umbrales de alerta con base en datos de referencia. La experiencia que se incorpora debe etiquetarse (por ejemplo, desafiante, reward escasa, cambio de entorno) para facilitar el análisis posterior.

\begin{importantbox}{Advertencia sobre la cobertura y la distribución del Replay Buffer}
Si el buffer se llena justamente con experiencias exitosas tempranas, las muestras de la exploración posterior no pueden entrar al entrenamiento, lo que provoca que el agent converja prematuramente. Se recomienda configurar \texttt{min\_size} y \texttt{max\_size}, revisar periódicamente las marcas de tiempo mediante la interfaz de muestreo y, cuando sea necesario, usar \texttt{prioritized replay} o \texttt{rank-based} para volver a incentivar el muestreo.
\end{importantbox}

\subsection{Estrategias de repetición de experiencia y rebalanceo de datos}

\begin{knowledgebox}{Técnicas de rebalanceo de Experience Replay}
\begin{itemize}
    \item Prioritized Replay: usa $|\delta|^\alpha$ como peso de muestreo, donde $\delta$ es el TD error y $\alpha$ controla la intensidad de la preferencia.
    \item Importance Sampling Correction: aplica una compensación proporcional a las muestras priorizadas para evitar el policy bias.
    \item Diversity Buffer: conserva fragmentos de experiencia provenientes de distintas subtareas o de distintas seed de simulación, para evitar que el entrenamiento quede atrapado en una narrow distribution.
\end{itemize}
\end{knowledgebox}

No basta con tratar el replay buffer como una caja negra; se recomienda registrar, después de cada iteración de entrenamiento, la proporción de muestras nuevas y antiguas que ingresan y su reward promedio. Si la proporción de muestras nuevas es menor que 30\%, la eficiencia de la exploración está disminuyendo y es necesario aumentar la tasa de aprendizaje o epsilon.

\begin{warningbox}{Dependencia excesiva del Replay buffer en datos antiguos}
Conforme avanza el entrenamiento, los datos antiguos pueden quedar obsoletos, sobre todo en entornos no estacionarios. Al continuar el entrenamiento hay que vigilar el timestep promedio del buffer; si más del 50\% de los datos proviene de los últimos diez minutos, se puede activar una limpieza cíclica o restablecer la prioridad a un valor alto, para evitar que la política antigua domine el gradiente.
\end{warningbox}

\subsection{Resumen de la sección}

Esta sección enfatiza la observabilidad en la construcción de sistemas de Q-Learning y la salud del replay buffer: todo reentrenamiento a gran escala debe establecer guardrail en tres dimensiones: la recolección de datos, la distribución de muestreo y el TD error.
